\DocumentMetadata{
  pdfversion=2.0,
  pdfstandard=ua-2,
  lang=en-US,
  tagging=on,
  tagging-setup={math/setup=mathml-SE,tabsorder=structure}
}
\documentclass[11pt,letterpaper]{article}
\usepackage[margin=0.68in,includefoot]{geometry}
\usepackage{newcomputermodern}
\usepackage{amsmath,amscd,mathtools}
\providecommand{\square}{\mdlgwhtsquare}
\usepackage[hidelinks]{hyperref}
\hypersetup{
  pdftitle={Complex Analysis PhD Exam, August 2007},
  pdfauthor={University of Florida Department of Mathematics},
  pdfsubject={Graduate mathematics examination},
  pdfdisplaydoctitle=true,
  bookmarksopen=true
}
\setcounter{secnumdepth}{0}
\setlength{\parindent}{0pt}
\setlength{\parskip}{0.28em}
\setlength{\emergencystretch}{3em}
\setlength{\leftmargini}{2.15em}
\setlength{\leftmarginii}{2.65em}
\setlength{\leftmarginiii}{3em}
\renewcommand{\labelenumi}{\textbf{\theenumi.}}
\pagestyle{empty}
\raggedbottom
\allowdisplaybreaks
\newenvironment{examproblems}[1]{%
  \begin{enumerate}%
  \setcounter{enumi}{#1}%
  \setlength{\topsep}{0.35em}%
  \setlength{\partopsep}{0pt}%
  \setlength{\itemsep}{0.48em}%
  \setlength{\parsep}{0.18em}%
}{\end{enumerate}}
\newenvironment{examsubparts}{%
  \begin{enumerate}%
  \setlength{\topsep}{0.18em}%
  \setlength{\partopsep}{0pt}%
  \setlength{\itemsep}{0.16em}%
  \setlength{\parsep}{0.1em}%
}{\end{enumerate}}
\newenvironment{examinstructions}{%
  \par\begingroup\itshape
}{%
  \par\endgroup\vspace{0.18em}
}
\makeatletter
\renewcommand\section{\@startsection{section}{1}{0pt}%
  {-1.25ex plus -.25ex minus -.1ex}%
  {0.55ex plus .1ex}%
  {\normalfont\large\bfseries}}
\renewcommand{\@maketitle}{%
  \newpage\null\vskip -1.2em%
  {\footnotesize\bfseries
    \href{https://www.ufl.edu/}{University of Florida}, \href{https://math.ufl.edu/}{Department of Mathematics}\par}%
  \vskip 0.18em%
  \tagstructbegin{tag=Title}%
    {\LARGE\bfseries\href{https://gma.math.ufl.edu/exams/complex-analysis-phd/}{Complex Analysis PhD Exam}, August 2007\par}%
  \tagstructend%
  \vskip 0.35em%
  \tagmcbegin{artifact}%
    \hrule height 1pt%
  \tagmcend%
  \vskip 0.55em%
}
\makeatother
\title{Complex Analysis PhD Exam, August 2007}
\author{}
\date{}
\begin{document}
\maketitle
\thispagestyle{empty}
\begin{examproblems}{0}
\item
Let \(\overline D\) be the closure of the unit disk. Assume \(f:\overline D\to\mathbb C\) is continuous with~\(f\) analytic on~\(D\), \(|f(z)|>3\) for \(|z|=1\), and \(f(0)=1-2i\). Must~\(f\) have a zero in the unit disk? Prove your answer is correct.
\item
Let~\(G\) be a simply connected region. A function \(f:G\to G\) is said to be biholomorphic if~\(f\) is bijective, analytic, and \(f^{-1}\) is also analytic. If \(z_1,z_2\) are distinct elements of~\(G\) and \(f_1,f_2:G\to G\) are biholomorphisms with \(f_1(z_1)=f_2(z_1)\) and \(f_1(z_2)=f_2(z_2)\), prove that \(f_1=f_2\).
\item
Assume that~\(f\) is entire, \(f(0)=3+4i\), and \(|f(z)|\leq 5\) when \(|z|<1\). What is \(f'(0)\)?
\item
Let \(\mathcal F\) be the collection of analytic mappings of open unit disk~\(D\) into \(\{\operatorname{Re}(z)>0\}\) such that \(f(0)=1\). Show that \(\mathcal F\) is a normal family.
\item
Let \(h_n\) be a sequence of harmonic functions on the connected open set~\(G\) and assume that \(h_n\to h\) uniformly on compact subsets of~\(G\). Show that~\(h\) is harmonic.
\item
Assume that~\(f\) and~\(g\) are entire and for all \(z\in\mathbb C\), \(|f(z)|\leq |g(z)|\). Show that for some constant~\(c\) with \(|c|\leq 1\), \(f=cg\).
\item
Assume that~\(f\) is a meromorphic function on \(\mathbb C\). A complex number~\(w\) is called a period of~\(f\) if \(f(z+w)=f(z)\) for all \(z\in\mathbb C\).
\begin{examsubparts}
\item[\textnormal{(a)}]
If \(w_1\) and \(w_2\) are periods of~\(f\), show that for all integers \(n_1\) and \(n_2\), \(n_1w_1+n_2w_2\) is also a period of~\(f\).
\item[\textnormal{(b)}]
If~\(G\) is a bounded region of \(\mathbb C\), show that~\(f\) has at most finitely many periods in~\(G\).
\end{examsubparts}
\item
Let~\(p\) be the polynomial \(p(z)=a_0+a_1z+\cdots+a_nz^n\). Show that for each \(j=0,1,\ldots,n\), 
\[
|a_j|\leq\max\{|p(z)|:|z|=1\}. \tag{1}
\]
\end{examproblems}
\end{document}
